%============================
% \section{ }
% \subsection{ }\label{ }\index{ }
%============================

\section{Markdown - Spickzettel}\label{markdown-syntax}

\subsection{Überschrift}\label{ueberschrift}

% Quellcode
\lstset{language=Bash} % C, [LaTeX]TeX, Bash, Python
\begin{lstlisting}[numbers=left, frame=l, framerule=0.1pt,%
% ======================
	caption={ }, % Caption
	label={code: }  % Label
]% =====================

    # Überschrift
    ## Überschrift 2
    ### Überschrift 3
\end{lstlisting}

\subsection{Bild}\label{bild}

% Bild Referenz
(\autoref{pic:logo1} logo.pdf). % Bildverweis = logo.pdf

\begin{figure}[!hp]% hier
	\centering
	\includegraphics[width=0.8\textwidth]{img/logo.pdf}
	% ====================
	\caption[logo.pdf ]{logo.pdf } % Caption
	\label{pic:logo1}	% Referenz
	% ====================
\end{figure}

% Quellcode
\lstset{language=Bash} % C, [LaTeX]TeX, Bash, Python
\begin{lstlisting}[numbers=left, frame=l, framerule=0.1pt,%
% ======================
	caption={ }, % Caption
	label={code: }  % Label
]% =====================

    # bild
    ![Logo](img/logo.pdf)
\end{lstlisting}

\subsection{Tabelle}\label{tabelle}

(\autoref{tab:Tabelle } ). % Tabellenverweis = Tabelle

% Tabelle
\begin{table}[!hp] % hier
	\centering
	%\setlength{\tabcolsep}{5mm} % Spaltenlänge fest
	\rowcolors{1}{}{lightgray!20}  % Farbe
	%\begin{tabularx}{\textwidth}{XX} % auto. Spaltenumbruch
	\begin{tabular} []{rll}
	\toprule % Inhalt
	% ---------
\textbf{Nr.} & \textbf{Begriffe} & \textbf{Erklärung}\\
	\midrule
1 & a1 & a2\\
2 & b1 & b2\\
3 & c1 & c2\\
4 & a1 & a2\\
	\bottomrule
	%\end{tabularx}
	\end{tabular}
	\caption[Tabelle ]{Tabelle }	% Caption
	\label{tab:Tabelle }	% Referenz
	%============================
\end{table}

% Quellcode
\lstset{language=Bash} % C, [LaTeX]TeX, Bash, Python
\begin{lstlisting}[numbers=left, frame=l, framerule=0.1pt,%
% ======================
	caption={ }, % Caption
	label={code: }  % Label
]% =====================

    # tabelle
    |**Nr.**|**Begriffe**|**Erklärung**|
    |--:|:-----|:----|
    | 1     | a1         | a2                |
    | 2     | b1         | b2              |
    | 3     | c1         | c2              |
    | 4     | a1         | a2              |
\end{lstlisting}

\subsection{Mathe}\label{mathe}

$[ V ] = [ \Omega ] \cdot [ A ]$ o. $U = R \cdot I$ o.
$R = \frac{U}{I}$

% Quellcode
\lstset{language=Bash} % C, [LaTeX]TeX, Bash, Python
\begin{lstlisting}[numbers=left, frame=l, framerule=0.1pt,%
% ======================
	caption={ }, % Caption
	label={code: }  % Label
]% =====================

    # Mathe
    $[ V ] = [ \Omega ] \cdot [ A ]$ o. $U = R \cdot I$ o. $R = \frac{U}{I}$
\end{lstlisting}

\textbf{Matheumgebung:}

\begin{align*}
    \sum_{i=1}^5 a_i = a_1 + a_2 + a_3 + a_4 + a_5
\end{align*}

% Quellcode
\lstset{language=Bash} % C, [LaTeX]TeX, Bash, Python
\begin{lstlisting}[numbers=left, frame=l, framerule=0.1pt,%
% ======================
	caption={ }, % Caption
	label={code: }  % Label
]% =====================

    # Matheumgebung
    \begin{align*}
        \sum_{i=1}^5 a_i = a_1 + a_2 + a_3 + a_4 + a_5
    \end{align*}
\end{lstlisting}

\subsection{Absätze}\label{absaetze}

Dies hier ist ein Blindtext zum Testen von Textausgaben. Wer diesen Text
liest, ist selbst schuld. Der Text gibt lediglich den Grauwert der
Schrift an. Ist das wirklich so? Ist es gleichgültig, ob ich schreibe:
\frqq Dies ist ein Blindtext\flqq\, oder \frqq Huardest gefburn\flqq\,? Kjift -
mitnichten! Ein Blindtext bietet mir wichtige Informationen.

Dies hier ist ein Blindtext zum Testen von Textausgaben. Wer diesen Text
liest, ist selbst schuld. Der Text gibt lediglich den Grauwert der
Schrift an. Ist das wirklich so? Ist es gleichgültig, ob ich schreibe:
\frqq Dies ist ein Blindtext\flqq\, oder \frqq Huardest gefburn\flqq\,? Kjift -
mitnichten! Ein Blindtext bietet mir wichtige Informationen.

\subsection{Listen}\label{listen}

\textbf{ungeordnete Liste}

\begin{itemize}% Liste Punkt
\item
  a
\item
  b

  \begin{itemize}% Liste Punkt
  \item
    bb
  \end{itemize}
\item
  c
\end{itemize}

% Quellcode
\lstset{language=Bash} % C, [LaTeX]TeX, Bash, Python
\begin{lstlisting}[numbers=left, frame=l, framerule=0.1pt,%
% ======================
	caption={ }, % Caption
	label={code: }  % Label
]% =====================

    # ungeordnete Liste
    - a
    - b
        - bb
    - c
\end{lstlisting}

\textbf{Sortierte Liste}

\begin{enumerate}% Liste 1) oder a)
\item
  eins
\item
  zwei
\item
  drei
\end{enumerate}

% Quellcode
\lstset{language=Bash} % C, [LaTeX]TeX, Bash, Python
\begin{lstlisting}[numbers=left, frame=l, framerule=0.1pt,%
% ======================
	caption={ }, % Caption
	label={code: }  % Label
]% =====================

    # Sortierte Liste
    1. eins
    2. zwei
    3. drei
\end{lstlisting}

\textbf{Sortierte Liste}

\begin{enumerate}% Liste 1) oder a)
\def\labelenumi{\alph{enumi})}
\item
  a
\item
  b
\item
  c
\end{enumerate}

(\autoref{code:Liste } Liste). % Codeverweis = Liste

% Quellcode
\lstset{language=Bash} % C, [LaTeX]TeX, Bash, Python
\begin{lstlisting}[numbers=left, frame=l, framerule=0.1pt,%
% ======================
	caption={Liste }, % Caption
	label={code:Liste }  % Label
]% =====================

    # Sortierte Liste
    a) a
    b) b
    c) c
\end{lstlisting}

\subsection{Anführungszeichen}\label{anfuehrungszeichen}

\frqq Anführungszeichen\flqq\, oder \flqq Anführungszeichen\frqq   oder
\frqq Anführungszeichen\flqq

% Quellcode
\lstset{language=Bash} % C, [LaTeX]TeX, Bash, Python
\begin{lstlisting}[numbers=left, frame=l, framerule=0.1pt,%
% ======================
	caption={ }, % Caption
	label={code: }  % Label
]% =====================

    # Anführungszeichen
    "Anführungszeichen" oder
    \flqq Anführungszeichen\frqq   oder
    \frqq Anführungszeichen\flqq
\end{lstlisting}

\subsection{Quellenangabe}\label{quellenangabe}

Quelle \textcite{schlosser_latex:2016}

Quelle Text\footfullcite{schlosser_latex:2016}

Quelle \cite{schlosser_latex:2016}

% Quellcode
\lstset{language=Bash} % C, [LaTeX]TeX, Bash, Python
\begin{lstlisting}[numbers=left, frame=l, framerule=0.1pt,%
% ======================
	caption={ }, % Caption
	label={code: }  % Label
]% =====================

    # Quellenangabe
    Quelle \textcite{schlosser_latex:2016}
    Quelle Text\footfullcite{schlosser_latex:2016}
    Quelle \cite{schlosser_latex:2016}
\end{lstlisting}
%============================
%============================
%============================
